\section{Vì sao phải chia cho căn của \texorpdfstring{\(\dk\)}{d k}}
\label{sec:ch07-chia-can}

\index{scaled dot-product!vì sao là tích vô hướng}%
Hộp cầu nối ở mục~\ref{sec:ch06-cau-noi} dừng lại ở một chuyện chưa xong. Luong
và cộng sự thử bốn hàm chấm điểm và không làm cho \code{concat}, tức họ hàng gần
nhất của hàm Bahdanau, chạy tốt được. Bài này chọn hàm rẻ nhất trong cả nhóm,
là \tn{tích vô hướng}{dot product}, và mục 3.2.1 nói rõ vì sao:

\begin{quote}
\enquote{Dot-product attention is identical to our algorithm, except for the
scaling factor of \(1/\sqrt{d_k}\). Additive attention computes the
compatibility function using a feed-forward network with a single hidden layer.
While the two are similar in theoretical complexity, dot-product attention is
much faster and more space-efficient in practice, since it can be implemented
using highly optimized matrix multiplication code.}
\end{quote}

Đọc câu ấy cạnh \eqref{eq:ch06-ham-cham-diem} thì thấy chỗ đổi. Hàm của Bahdanau
cộng hai phép chiếu rồi cho qua tanh rồi chiếu xuống một số thực; hàm ở đây nhân
hai vector với nhau. Bài nói hai bên \enquote{similar in theoretical
complexity} và chỗ khác nhau nằm ở chữ \enquote{in practice}: một phép chấm
điểm viết được thành một phép nhân ma trận thì chạy trên đúng đoạn mã mà mọi
thư viện đại số tuyến tính đã tối ưu sẵn, còn một mạng một tầng ẩn thì không.

Phương trình (1) của bài, gọn hết mức:

\begin{equation}
  \attention(\mat{Q}, \mat{K}, \mat{V})
  = \softmax\!\left( \frac{\mat{Q}\mat{K}\transpose}{\sqrt{\dk}} \right) \mat{V}.
  \label{eq:ch07-attention}
\end{equation}

Trong repo companion, dòng ấy là bốn dòng:

\begin{minted}{python}
d_k = query.shape[-1]
scores = query @ key.transpose(-2, -1) / math.sqrt(d_k)
if mask is not None:
    scores = scores.masked_fill(~mask, float("-inf"))
weights = torch.softmax(scores, dim=-1)
\end{minted}

Cái đáng dừng lại là mẫu số. Bài đặt tên cho cả cơ chế theo nó, scaled
dot-product attention, nên nó không phải một hằng số dọn dẹp.

\subsection*{Dẫn xuất}

\index{scaled dot-product!phương sai của tích vô hướng}%
Chú thích số 4 của bài đưa lập luận trong hai câu:

\begin{quote}
\enquote{To illustrate why the dot products get large, assume that the
components of \(q\) and \(k\) are independent random variables with mean 0 and
variance 1. Then their dot product, \(q \cdot k = \sum_{i=1}^{d_k} q_i k_i\),
has mean 0 and variance \(d_k\).}
\end{quote}

Dẫn ra thì mất ba dòng. Gọi \(q_i\) và \(k_i\) là các biến ngẫu nhiên độc lập,
kỳ vọng 0, phương sai 1. Với mỗi số hạng,

\begin{equation}
  \E[q_i k_i] = \E[q_i]\,\E[k_i] = 0,
  \qquad
  \E[(q_i k_i)^2] = \E[q_i^2]\,\E[k_i^2] = 1,
  \label{eq:ch07-so-hang}
\end{equation}

nên mỗi số hạng có phương sai 1. Các số hạng độc lập nhau, và phương sai của
tổng các biến độc lập là tổng các phương sai, nên

\begin{equation}
  \E[\vect{q} \cdot \vect{k}] = 0,
  \qquad
  \operatorname{Var}[\vect{q} \cdot \vect{k}] = \dk .
  \label{eq:ch07-phuong-sai}
\end{equation}

Độ lệch chuẩn là \(\sqrt{\dk}\). Nên điểm số vào softmax không có thang cố
định: nó nở ra theo căn của bề rộng một head. Ở \(\dk = 64\) của bài, điểm số điển
hình cỡ 8; ở \(\dk = 512\) cỡ 22. Chia cho \(\sqrt{\dk}\) đưa độ lệch chuẩn về
1 và giữ nguyên ở đó.

\subsection*{Đo, vì lập luận trên giấy không nói cái giá}

\index{scaled dot-product!đo phương sai}%
Dẫn xuất nói điểm số lớn lên. Nó không nói lớn lên thì hỏng ra sao, và bài chỉ
nói bằng một mệnh đề: softmax bị đẩy vào vùng \enquote{where it has extremely
small gradients}. Đó là một phát biểu đo được, nên đo.

Trước hết kiểm chính chú thích ấy.

\begin{measured}{phương sai của tích vô hướng theo \(\dk\), \(\vect{q}, \vect{k}\) chuẩn}
\begin{minted}{text}
d_k    mean       var         var/d_k   sd        sqrt(d_k)
8      0.0226     8.0037      1.0005    2.8291    2.8284
16     0.0594     15.9820     0.9989    3.9977    4.0000
32     0.0047     32.1030     1.0032    5.6659    5.6569
64     0.0533     63.7511     0.9961    7.9844    8.0000
128    0.0184     125.6465    0.9816    11.2092   11.3137
256    -0.0128    256.9861    1.0039    16.0308   16.0000
512    0.1007     508.2525    0.9927    22.5445   22.6274
1024   -0.2621    1034.2814   1.0100    32.1602   32.0000
\end{minted}

20000 mẫu mỗi dòng.

Đo bằng \code{experiments/ch07\_scaling.py} tại tag \code{ch07}.
\end{measured}

Cột \code{var/d\_k} nằm giữa 0.9816 và 1.0100 trên tám bậc gấp đôi. Chú thích
đúng.

Bây giờ tới cái giá. Dựng một hàng attention trên 64 khóa, chấm điểm, rồi đọc
ba thứ: trọng số lớn nhất, entropy của phân phối, và chuẩn Frobenius của
\tn{ma trận Jacobi}{Jacobian matrix} của softmax. Cái thứ ba là con số dịch cụm \enquote{extremely small
gradients} sang tiếng toán: mọi gradient đi ngược qua softmax đều nhân với ma
trận ấy, nên chuẩn của nó tiến về 0 là gradient tiến về 0.

\begin{measured}{softmax trên 64 khóa, chia và không chia}
\begin{minted}{text}
       unscaled                    scaled by 1/sqrt(d_k)
d_k    max_p   entropy  grad       max_p   entropy  grad
8      0.4370  1.9930   0.303610   0.1081  3.6715   0.184700
16     0.6137  1.2680   0.281309   0.1156  3.6583   0.187285
32     0.7398  0.7896   0.235638   0.1120  3.6762   0.184970
64     0.7795  0.6256   0.219978   0.1034  3.6929   0.182454
128    0.8542  0.3930   0.165504   0.1034  3.7018   0.181279
256    0.9055  0.2494   0.114576   0.1061  3.6906   0.182846
512    0.9263  0.1844   0.091270   0.1103  3.6749   0.185140
1024   0.9431  0.1429   0.077089   0.1064  3.6870   0.183148
\end{minted}

Trung bình trên 200 hàng. \code{grad} là chuẩn Frobenius của Jacobi softmax.

Đo bằng \code{experiments/ch07\_scaling.py} tại tag \code{ch07}.
\end{measured}

\index{scaled dot-product!chỗ phép chia làm gì}%
Ba cột bên trái là chuyện bài mô tả. Trọng số lớn nhất đi từ 0.4370 lên 0.9431,
tức phân phối dồn dần về một cột; entropy tụt từ 1.9930 xuống 0.1429; và
gradient tụt từ 0.303610 xuống 0.077089, gần bốn lần.

Ba cột bên phải là chỗ tôi muốn bạn dừng lại, và nó không phải chuyện
\enquote{chia thì gradient lớn hơn}. Ở \(\dk = 8\) phép chia làm gradient
\emph{nhỏ} đi, 0.184700 so với 0.303610. Cái nó làm là khác: entropy nằm giữa
3.6583 và 3.7018, gradient nằm giữa 0.181279 và 0.187285, trên toàn bộ dải
\(\dk\) từ 8 tới 1024.

Phát biểu đúng là thế này. Phép chia không làm softmax nhạy hơn. Nó làm softmax
\emph{hết phụ thuộc vào \(\dk\)}. Bạn đổi bề rộng một head từ 8 lên 1024 và hành vi
của lớp softmax đứng yên, nên chọn \(\dk\) trở thành một quyết định về sức chứa
chứ không còn là một quyết định về độ \tn{bão hòa}{saturation}. Đó cũng là chỗ
nối với hàng (A)
của bảng 3, mà mục~\ref{sec:ch07-nhieu-dau} đọc: bảng ấy đổi số head, tức đổi
\(\dk\), và nếu không có phép chia thì mỗi hàng của nó còn đang đo thêm một
biến nữa.

\subsection*{Giả thiết của chú thích không phải cái mã thật làm}

\index{scaled dot-product!giả thiết N(0,1) không đúng với mã thật}%
Còn một chuyện nữa, và tôi ghi ra vì nó là chỗ giấy và máy lệch nhau.

Chú thích giả định các thành phần của \(\vect{q}\) và \(\vect{k}\) có phương sai
1. Nhưng \(\vect{q}\) không rơi từ trời xuống: nó là \(\mat{W}_Q \vect{x}\), và
\(\mat{W}_Q\) khởi tạo theo mặc định của framework. Mặc định của
\code{nn.Linear} trong PyTorch là phân phối đều trên
\([-1/\sqrt{n_{\mathrm{in}}},\, 1/\sqrt{n_{\mathrm{in}}}]\), không phải chuẩn
đơn vị.

\begin{measured}{phương sai thật của một phép chiếu ở lúc khởi tạo}
\begin{minted}{text}
d_model  d_k    var(component)  var(q.k)   d_k     ratio
512      64     0.333680        7.174865   64      0.112107
512      512    0.333187        56.362167  512     0.110082
64       16     0.328320        1.692329   16      0.105771
64       64     0.330210        7.071695   64      0.110495
\end{minted}

Đầu vào chuẩn đơn vị, \code{Linear} không bias, khởi tạo mặc định.

Đo bằng \code{experiments/ch07\_scaling.py} tại tag \code{ch07}.
\end{measured}

Phương sai mỗi thành phần là 1/3 chứ không phải 1, và con số ấy dẫn ra được:
phương sai của phân phối đều trên \([-b, b]\) là \(b^2/3\), với
\(b = 1/\sqrt{n_{\mathrm{in}}}\) thì tổng qua \(n_{\mathrm{in}}\) chiều cho
đúng 1/3. Kéo theo tích vô hướng có phương sai \(\dk/9\) chứ không phải
\(\dk\), và cột \code{ratio} đo được 0.112107, 0.110082, 0.105771, 0.110495,
tức quanh \(1/9\).

Nghĩa là ở lúc khởi tạo, \(\sqrt{\dk}\) chia quá tay khoảng ba lần.

Chuyện đó \emph{không} làm hỏng lập luận, và đây là chỗ đáng đọc chậm. Việc của
phép chia là bỏ sự phụ thuộc vào \(\dk\), không phải đặt một thang tuyệt đối.
Một hằng số nhân dư ba lần là thứ \(\mat{W}_Q\) và \(\mat{W}_K\) hấp thụ được,
vì chúng là tham số học được và một phép nhân vô hướng nằm gọn trong không gian
chúng đi tới; còn một hệ số nở theo \(\sqrt{\dk}\) thì không hấp thụ được như
thế, vì nó đổi theo một siêu tham số mà tối ưu hóa không nhìn thấy. Bảng thứ
hai ở trên đo đúng chuyện ấy: sau khi chia, ba cột đứng yên theo \(\dk\).

Tôi không đo mất bao nhiêu bước để phần dư ba lần ấy bị hấp thụ, và cũng không
tìm thấy ai đo. Câu ở trên là câu về chuyện \emph{có thể} hấp thụ được, không
phải câu về tốc độ.

Tôi ghi lại vì nó là một thói quen đọc paper hơn là một chi tiết về PyTorch. Một
dẫn xuất có giả thiết, và giả thiết ấy hiếm khi đúng nguyên văn trong mã. Câu
hỏi phải hỏi không phải \enquote{giả thiết có đúng không} mà là \enquote{kết
luận còn đứng không khi giả thiết sai lệch theo kiểu này}. Ở đây còn, và lý do
là kết luận nói về một sự phụ thuộc chứ không nói về một con số.

\textbf{Chưa đo}: thống kê ấy sau khi huấn luyện xong. Mọi số trong mục này là
số ở lúc khởi tạo, và bài tập cuối chương giao lại phép đo kia.
